\section{System overview and design principles}
\label{sec:system}

QDCVR (v2.3) is a self-hosted knowledge-base management system for converting
heterogeneous documents into searchable, governable collections. A Python
parsing service, a TypeScript web application and an MCP tool server share
local storage and retrieval engines \cite{mcp}.

\subsection{Design principles: three contracts and enforcement layers}
\label{sec:req}
\label{sec:principles}

The three contracts are implemented at different layers. Some checks are
mechanically enforced by server-side tools; placement and answer abstention
also depend on the calling workflow. We therefore distinguish the mechanism
that executes each step from the audit or evaluation that checks its outcome.

\begin{description}
\item[Content informs placement.] \emph{Mechanism:} an ingest agent reads the
parsed content together with the library catalogue and selects a parent or
sub-library; storage APIs validate identifiers and required fields, not the
semantic correctness of that choice. \emph{Check:} descriptions and tags are
re-read against the body, and three-source consistency is checked after
lifecycle operations (Sections~\ref{sec:acquisition} and~\ref{sec:organization}).
\item[Content informs admission.] \emph{Mechanism:} when invoked, the server-side
Laya judging tool fetches or accepts document text, scores its segments, and
returns survivors with provenance; lane workflows are instructed to treat rank
as advisory. \emph{Check:} failed fetches or scores are not admitted, and the
ablation in Section~\ref{sec:e8} tests a bounded frozen-pool contrast
(Sections~\ref{sec:lanes} and~\ref{sec:gate}).
\item[Absence is reported.] \emph{Mechanism:} the answer workflow instructs the
agent to report not-found when completed retrieval provides no supporting
evidence and to state its searched scope. This is not a server-side semantic
classifier. \emph{Check:} out-of-corpus probes and unanswerable questions were
observed as abstentions in the reported trials (Section~\ref{sec:e5}).
\end{description}

Four deployment requirements constrain the contracts: expose parsing failures;
keep document files and library records consistent; require content evidence
before a retrieved item is used to support an answer; and remain operable on
one workstation. Server-side checks cover parse, metadata and index
post-conditions, the invoked evidence gate, and endpoint tool policies. They do
not mechanically establish semantic placement accuracy or answer faithfulness.

\subsection{Architecture and subsystems}
\label{sec:architecture}

Fig.~\ref{fig:arch} shows the acquisition and query paths. Ingestion parses and
quality-checks documents; an agent uses the parsed content and library
catalogue to select a parent or sub-library; save and index operations then
update records and search indices. At query time, the chat endpoint supplies a
lane-specific tool policy, and the agent retrieves candidates and invokes the
gate before the answer layer returns cited evidence or not-found. The MCP
server exposes 95 registered tools. The three contracts of
Section~\ref{sec:principles} sit at different points along this path:
placement during ingestion, admission at the gate, the abstention report in
the answer layer.

The management path covers acquisition (Section~\ref{sec:acquisition}),
content-routed organization and lifecycle checks
(Section~\ref{sec:organization}), and the tool and runtime layer
(Section~\ref{sec:implementation}). The chat
endpoint defines pinned-search and all-library librarian policies. Hybrid is
a composed retrieval configuration, not a third built-in server policy; it
combines vector and catalog calls under an explicit directive. The document
graph supports navigation, not evidence, and the experience store uses the
same catalog and gate but was not evaluated at scale.

\begin{figure*}[t]
\centering
\includegraphics[width=\textwidth]{fig1_architecture}
\caption{QDCVR architecture and three design commitments. During ingestion, an
agent uses document content and the catalogue to choose a library; server-side
storage and consistency checks verify structural state (C1). At query time,
the endpoint applies a tool policy and the agent may invoke the server-side
Laya evidence tool, which returns scored survivors with provenance (C2). The
answer workflow is instructed to report not-found after a completed search
finds no support (C3). These commitments span distinct enforcement layers;
the evaluation measures selected mechanisms and observed outcomes.}
\label{fig:arch}
\end{figure*}

\subsection{Data model and object lifecycle}
\label{sec:datamodel}
\label{sec:lifecycle}

Documents are stored as Markdown with extracted images. Three sources describe
library state: document files, the \texttt{.tree-fs.json} registry and each
library's \texttt{.knowledge-base.yml}. Ingestion proceeds through raw,
parsed, stored and indexed states, and is complete only after a content-derived
probe retrieves the document. Indexing is explicit and asynchronous; task
identifiers support polling, while post-operation audits detect stale indices
or disagreement among the three sources.
